% ============================================
% Senior Associate Cover Letter – Vega Economics
% ============================================

\documentclass[11pt,letterpaper]{letter}

\usepackage{geometry}
\usepackage{microtype}
\usepackage[hidelinks]{hyperref}

\geometry{left=25mm,right=25mm,top=25mm,bottom=25mm}

\signature{Decory Edwards}
\address{
Department of Economics \\
Johns Hopkins University \\
Baltimore, MD 21218 \\
\texttt{dedwar65@jh.edu}
}

\begin{document}

\begin{letter}{Hiring Committee \\
Vega Economics \\
Berkeley, CA \\ USA}

\opening{Dear Members of the Hiring Committee,}

I am writing to apply for the Senior Associate position at Vega Economics. I am a Ph.D.\ candidate in Economics at Johns Hopkins University, advised by Professors Christopher D.~Carroll, Nicholas W.~Papageorge, and Stelios Fourakis, and I expect to complete my degree in May~2026. My research fields are macroeconomics, wealth and income dynamics, and computational economics.

My research agenda focuses on understanding the determinants of wealth inequality and the mechanisms through which heterogeneity in returns to wealth shapes aggregate outcomes. In my job-market paper, I estimate the degree of heterogeneity in individual portfolio returns using micro-data from the Survey of Consumer Finances and simulate a structural model in which households face idiosyncratic return risk. I show that allowing for persistent heterogeneity in returns substantially improves the model’s ability to match the empirical distribution of wealth, offering a tractable way to connect micro-level return dispersion to macroeconomic inequality.

In a second project, I use the Health and Retirement Study to examine the relationship between interpersonal trust and portfolio performance. Building on the literature that finds a hump-shaped relationship between trust and income, I show that a similar relationship holds for individual returns to wealth measured in the HRS data, highlighting a behavioral channel through which social attitudes shape saving and investment outcomes. This work also provides new estimates of the persistent component of returns to net worth, which motivate the calibration of heterogeneity in my structural model.

Vega Economics’ focus on competition and antitrust matters is a strong fit for my quantitative training and applied interests. Your work requires translating economic theory and empirical evidence into clear, defensible conclusions. This includes evaluating competitive effects, market power, and counterfactual outcomes using detailed market and firm data, and constructing damages analyses that are transparent, replicable, and robust. My research training has prepared me to build and assess counterfactuals, implement empirical designs carefully, and communicate results clearly for technical and non-technical audiences.

In addition to research, I have extensive experience with data analysis in Python and Stata, producing reproducible workflows, and writing clearly under deadlines. I enjoy collaborative work, I am comfortable taking ownership of empirical workstreams, and I value precision and transparency in both code and written products, skills that are essential in litigation-oriented economic consulting.

I am particularly drawn to Vega Economics because of its reputation for rigorous analysis and close engagement on cases. I would be enthusiastic about contributing to your work in Berkeley and building as an applied economist in an environment where careful quantitative work directly informs high-stakes decisions.

Thank you very much for your time and consideration. I would welcome the opportunity to discuss my application further.

\closing{Sincerely,}

\end{letter}
\end{document}
